\section{Relative progression errors and simultaneous prime pairs}
\label{sec:simultaneous-prime-pairs}

This section continues CC1--CC20.  It supplies the analytic prime-selection
endpoints \texttt{PrimesRel}, \texttt{twistK}, \texttt{meanSqK},
\texttt{pair\_exists}, and \texttt{PairEx}, and the real-zero exclusion on
an already specified compatible host.  Constants denoted by $C,c,A$ are
positive and absolute unless their dependence is explicitly indicated.
Zeros are counted with analytic multiplicity.  No assertion about zeros
on the critical line is used here.

Three classical inputs are retained, with their mathematical content
specified here rather than hidden inside a source theorem name.
\begin{enumerate}
\item[(CN)] The analytic class-number formula for quadratic fields, with
the usual ideal class number $h$, number of roots of unity $w$, and
regulator $R$: the residue is
$2^{r_1}(2\pi)^{r_2}hR/(w\sqrt{|D|})$.
In particular a real primitive nonprincipal Dirichlet character of
conductor $f$ has $L(1,\chi)\ge c/\sqrt f$.
The convention for a real quadratic field is
$L(1,\chi_D)=2h\log\epsilon_D/\sqrt D$, with $\epsilon_D>1$
generating its units modulo torsion.
\item[(BD)] Bombieri's log-free density theorem, including its exceptional
factor.  There are absolute $A_B,C_B,c_B>0$ such that the total number
$N_R(\sigma)$ of zeros of primitive characters of conductor at most $R$,
with $\sigma\le\Re\rho<1$ and $|\Im\rho|\le R$, is at most
$C_BR^{A_B(1-\sigma)}$, for $R\ge2$ and $1/2\le\sigma\le1$.
If a real primitive character in this family has a real zero
$\beta=1-\delta>1-c_B/\log R$, omit that one simple zero; the resulting
count is at most
$C_B\delta\log R\,R^{A_B(1-\sigma)}$.
\item[(HD)] A classical hybrid density estimate: for each fixed $h>0$,
\[
 \sum_{f\le Q}\ \sum_{\chi\bmod f}^{*}N(\sigma,T,\chi)
       \le C_h(Q^2T)^{2(1-\sigma)+h},
 \quad Q\ge2, T\ge1, 11/14\le\sigma\le1.
\]
Here the asterisk restricts to primitive characters, and both signs of
the imaginary part are counted.  This is the form of Heath-Brown's
1979 Theorem 1 that will be used.
\end{enumerate}
These are named classical premises, not proofs reconstructed in this
section.  Primary references are Bombieri, \emph{Le grand crible dans la
th\'eorie analytique des nombres}, Ast\'erisque 18, \S6, Theorem 14
(the 1987 second edition preserves that statement), and Heath-Brown,
\emph{The density of zeros of Dirichlet's $L$-functions},
Canad. J. Math. 31 (1979), 231--240, Theorem 1,
\url{https://doi.org/10.4153/CJM-1979-024-0}.
For the normalization in (CN), see Sutherland's MIT 18.785 notes,
Lecture 19 (2021), on the analytic class-number formula.
The source's particular large numerical density constants and its
fixed-$1/8$ exceptional-pair wrapper are not asserted by these substitutions.

Write $\psi_\chi(u)=\sum_{n\le u}\Lambda(n)\chi(n)$ and
$\operatorname{Li}(u)=\int_2^u dt/\log t$.  Primitive principal characters
occur only at conductor one.  The tail estimate is first proved for
nonprincipal primitive characters and then extended to every nonprincipal
character at its original level by an explicit deleted-Euler-factor bound.

\paragraph{QP1: an effective real-zero gap with the right normalization.}
Suppose $f\le m$, $m\ge3$, and $\chi$ is a real primitive nonprincipal
character of conductor $f$.  The two quadratic cases of (CN) give the
stated uniform lower bound: for $D<0$, $w\le6$ and $h\ge1$; for $D>0$,
$h\ge1$ and $\epsilon_D\ge(1+\sqrt5)/2$.  The latter inequality follows
by taking the integer trace of a quadratic unit of norm $1$ or $-1$.

For $1-1/\log m\le s\le1$, periodic cancellation gives
$|\sum_{n\le u}\chi(n)|\le f$.  Split $L(s,\chi)$ at $m$ and differentiate
the absolutely convergent partial-summation integral for its tail.
The finite part is bounded by
$e\sum_{n\le m}(\log n)/n\le C\log^2m$.
The tail and its derivative are bounded by
$Cf m^{-s}(1+\log m)\le C(1+\log m)$, after decreasing the fixed
neighborhood of one if necessary.  Thus $|L'(s,\chi)|\le C\log^2m$.
If $L(\beta,\chi)=0$ in this neighborhood, the real mean-value theorem
implies
\[
 1-\beta\ge \frac{c}{\sqrt f\log^2m}
       \ge\frac{c'}{m^2\log m}\ge\frac{c'}{m^3}.       \tag{QP.1}
\]
The second inequality is uniform for $m\ge3$, since
$\log m\le m^{3/2}$.  The constants do not depend on the character.
This proves the effective lower-gap clause needed by
\texttt{PrimesRel}; it does not assert Siegel's ineffective lower bound.

\paragraph{QP2: one exceptional primitive character across bounded conductors.}
There is $c_P>0$ such that among primitive conductors at most $R$, at most
one real character has a zero $\beta>1-c_P/\log R$; that zero is simple.
Here is the positivity argument, including the across-conductor point.
For two distinct real primitive characters $\chi_1,\chi_2$, inflate them
to a common modulus.  The coefficients of
\[
 -\frac d{ds}\log\{\zeta(s)L(s,\chi_1)L(s,\chi_2)
                                      L(s,\chi_1\chi_2)\}
\]
are $\Lambda(n)(1+\chi_1(n))(1+\chi_2(n))\ge0$ on $s>1$.
The product character is nonprincipal: otherwise the primitive characters
would agree.  Removing Euler factors at the common modulus changes the
local logarithmic-derivative estimates by $O(\log R)$.
The real part of every critical-zero term at a real $s>1$ is positive.
The local expansion of CC3, or its Hadamard-product limit, therefore gives
\[
 \frac1{s-\beta_1}+\frac1{s-\beta_2}
                  \le\frac1{s-1}+C\log R.             \tag{QP.2}
\]
Set $s=1+a/\log R$.  First choose $a>0$ with $Ca<1/4$, and then
$c_P<a/10$.  The inequality
$2/(a+c_P)\le1/a+C$ is impossible.  Applying the same argument to
$\zeta(s)L(s,\chi)$, counting multiplicity, excludes two nearby real
zeros of one character and excludes a multiple nearby zero.
The standard trigonometric positivity argument of CC5 excludes the
nonreal alternatives in a sufficiently small region
$\Re s>1-c/\log(f(|\Im s|+3))$.
Consequently at a fixed modulus $m$ one can select a unique exceptional
pair, or select no pair, using a sufficiently small fixed $c_1\le1/8$.
Its primitive conductor divides $m$, and induction back to modulus $m$
is unique.  QP1 applies to this pair.

\paragraph{QP3: a weighted zero sum with the exceptional factor retained.}
Fix $m\ge3$, $T\ge2$, put $H=\log(m(T+3))$, and suppose
$V\ge C_0H$ for a sufficiently large absolute $C_0$.
For characters modulo $m$, omit the selected exceptional zero if present.
Their remaining zeros satisfy $\Re\rho\le1-c/H$ by QP2 and CC5.
Let $D=1$ when no exceptional pair is selected, and let
$D=\min(1,\delta H)$ when the selected zero is $1-\delta$.
Then (BD) gives, uniformly for $1/2\le\sigma\le1$,
\[
 N'(\sigma;m,T)\le C D\exp\{A_BH(1-\sigma)\}.          \tag{QP.3}
\]
Indeed use the primitive family with cutoff $R=m(T+3)$.  If
$\delta H<c_B$, the selected zero is exceptional in that larger family,
so the second part of (BD) applies.  If $\delta H\ge c_B$, its first part
suffices after multiplying by $\max(1,c_B^{-1})$.  Distinct characters
modulo $m$ have distinct primitive inducers.  Extra Euler factors have
zeros only on $\Re s=0$, outside the region being counted.

For $b=\Re\rho\ge1/2$, use
\[
 e^{-(1-b)V}=e^{-V/2}+V\int_{1/2}^{b}e^{-(1-\sigma)V}\,d\sigma.
\]
Summation, followed by (QP.3), bounds the sum by the two terms
\[
 CD e^{-(V-A_BH)/2},\qquad CDV\int_{c/H}^{1/2}e^{-(V-A_BH)t}\,dt .
\]
Take $C_0\ge2A_B+4$.  Both terms are at most
$CD e^{-c'V/H}$.  Since $|\rho|\ge1/2$, this proves
\[
 \sum_{\chi\bmod m}\ \sum_{\substack{\Re\rho\ge1/2, |\Im\rho|\le T\cr
                                    \rho\ {\rm not\ exceptional}}}
       \frac{e^{(\Re\rho-1)V}}{|\rho|}
             \le C D e^{-c'V/H}.                       \tag{QP.4}
\]
The factor $D$ is essential when the main term is small.  Replacing it
by $1$ would not prove a relative progression estimate.

\paragraph{QP4: the right-half explicit formula and relative Chebyshev error.}
Let $u=e^V$, $m\ge3$, and choose
$T=m^8e^{\sqrt V}$, with $u\ge m^{C_2}$ and $C_2$ sufficiently large.
Then $T\le u$ and $H\le10\log m+\sqrt V$, while $V\ge C_0H$.
Character orthogonality and the explicit formula yield
\[
 \varphi(m)\psi(u;m,a)
 =u-\theta\frac{u^\beta}{\beta}
   -\sum_{\chi\bmod m}\overline{\chi(a)}
      \sum_{\rho\ne\beta}\frac{u^\rho}{\rho}+E,         \tag{QP.5}
\]
where the displayed zero sum can be restricted to
$\Re\rho\ge1/2$, $|\Im\rho|\le T$, and
\[
 |E|\le C m\left\{u\log^2(muT)/T+
                              u^{1/2}\log^2(muT)\right\}.\tag{QP.6}
\]
For this right-half form there is no small-zero denominator left in the
error.  In the notation of CC8, the zero-removed local logarithmic
derivative gives the signed identity
\[
 b_0=\sum_{\rho\in Z_\chi(1)}\frac{m_\rho}{\rho}
                         +O(\log(N+2)).
\]
Indeed, each local term $-m_\rho/(s-\rho)$ has constant term
$m_\rho/\rho$ at zero; the omitted zeros at height at least one and
the regular gamma terms cost $O(\log(N+2))$ by CC3.
The right-half summands in this constant have $|\rho|\ge1/2$ and
are harmless.  For the left-half low zeros, combining the constant
with the explicit-formula sum gives exactly
$-m_\rho(u^\rho-1)/\rho$.  If $|\rho|$ is small and $\Re\rho\le1/2$, its
absolute value is at most $m_\rho u^{1/2}\log u$ using
$(u^\rho-1)/\rho=\int_1^u t^{\rho-1}dt$.
The CC3 unit-height counts then give the second term of (QP.6).
Zeros away from zero are bounded directly.  The removal of Euler factors
at primes dividing $m$ costs at most $C\log m\log u$, also absorbed there.
The CC8 floor-endpoint convention remains in force.

Here $\theta=0$ if no pair was selected; otherwise
$\theta=\chi_1(a)\in\{1,-1\}$ and $7/8\le\beta<1$.
For $V\ge8$, put $w=1$ in the former case and
$w=\min(1,(1-\beta)V)$ in the latter.  Elementary exponentiation gives
\[
 A(u):=u-\theta u^\beta/\beta\ge uw/4.                \tag{QP.7}
\]
For example $e^{-\delta V}/(1-\delta)\le e^{-\delta V/2}$
when $V\ge8$ and $\delta\le1/8$; this proves (QP.7) for the only delicate
case $\theta=1$.
Because $H\le V$, $D\le w$, so (QP.4) bounds the zero term relative to
$A(u)$ by $Ce^{-cV/H}$.
For the error (QP.6), (QP.1) gives $w\ge\min(1,c m^{-3}V)$.
Consequently its first relative part is at most
$C m^{-4}(V+H)^2e^{-\sqrt V}$, and its second is at most
$C m^4(V+H)^2e^{-V/2}$.  Increasing $C_2$ and decreasing $c$ absorbs
the polynomial factors uniformly.  Finally
$V/(10\log m+\sqrt V)\ge\frac1{20}\min(V/\log m,\sqrt V)$.
We have proved
\[
 |\varphi(m)\psi(u;m,a)-A(u)|
 \le C\{e^{-cV/\log m}+e^{-c\sqrt V}\}A(u).           \tag{QP.8}
\]
All constants are chosen before $m,a,u$ and before the exceptional pair.

\paragraph{QP5: partial summation preserves the relative error.}
Set $B(u)=\operatorname{Li}(u)-\theta\operatorname{Li}(u^\beta)$.
The identity $B'(u)=(1-\theta u^{\beta-1})/\log u$ exactly matches
$A'(u)/\log u$; in particular there is no extra factor $1/\beta$ in
the prime-counting main term.  Prime powers of exponent at least two
contribute $O(\sqrt u\log^2u)$ to $\psi$.
Apply partial summation on $[\sqrt u,u]$, using (QP.8) there after
increasing the starting power of $m$.  On this interval the error factor
is at most $C(e^{-cV/(2\log m)}+e^{-c\sqrt{V/2}})$.
The positive weight gives an error at most that factor times
\[
 A(u)/\log u+\int_{\sqrt u}^{u}A(t)/(t\log^2t)\,dt
                                  \le 2A(u)/\log u.
\]
Here $A(t)$ is increasing on the interval.  The portion below $\sqrt u$
and the prime-power correction cost $O(m\sqrt u\log^2u)$.

For $\theta=1$, integration over $[u^\beta,u]$ shows
$B(u)\ge c\,u\min(1,\delta\log u)/\log u$;
also $0<A(u)\le u\min(1,\delta\log u)$.
For $\theta=0$ or $-1$, both comparison statements hold with the minimum
replaced by $1$ and with changed absolute constants.  Thus the preceding
error is relative to $B(u)$, and (QP.1) absorbs the initial segment.
We obtain, for $u\ge m^{C_3}$,
\[
 |\varphi(m)\pi(u;m,a)-B(u)|
 \le C\{e^{-c\log u/\log m}+e^{-c\sqrt{\log u}}\}B(u).
                                                               \tag{QP.9}
\]
Taking $c_1$ from QP2, the zero-gap constant from QP1, and changing $c$
to match $e^{-\sqrt{c_3\log u}}$, gives the literal existential constants
of \texttt{GlobalTheory.Lin.PrimesRel} and its \texttt{Emain} convention.
This obtains that endpoint under (CN) and (BD).  It does not identify the
source's fixed-$1/8$ \texttt{PNTAPRelW} record; the argument below uses
(QP.9) directly and needs no such identification.

\paragraph{QP6: many primes in every admitted progression.}
Call $(Y,U,m,a)$ admitted at exponent $\epsilon>0$ when
$m>0$, $8\mid m$, $(a,m)=1$, and
\[
 \tfrac14\log^2Y\log m\le\log U,\qquad
 Y^{\epsilon/6}\le\log U\le Y^{\epsilon/5}\log^2Y.     \tag{QP.10}
\]
Let $P=\{q:\ q\text{ prime},\ \sqrt U<q\le U, q\equiv a\pmod m\}$.
For all sufficiently large $Y$, uniformly in the admitted data,
\[
                         |P|\ge U^{9/10}.              \tag{QP.11}
\]
Indeed (QP.10) gives $\log m=o(\log U)$ uniformly, so (QP.9) applies
and its relative error tends to zero.  Even in the exceptional case,
(QP.1) and QP5 give
$B(U)/\varphi(m)\ge cU/(m^4\log U)$.
Since $\log m\le4\log U/\log^2Y$ and
$\log\log U=O_\epsilon(\log Y)$, this lower bound is eventually
at least $2U^{9/10}$.  Removing at most $\sqrt U$ small primes proves
(QP.11).  This uses the effective lower gap, not an ineffective
constant depending on a moving modulus.

\paragraph{QP7: cancellation outside one squarefree exceptional index.}
For any fixed $A,B,\epsilon>0$ and all sufficiently large $Y$, admitted
data have an index $s_*\ge0$ such that for every odd squarefree
$1<s\le Y^A$, $(s,m)=1$, $s\ne s_*$,
\[
 \left|\sum_{q\in P}\left(\frac qs\right)\right|
                                        \le Y^{-B}|P|.\tag{QP.12}
\]
Put $R=mY^A$.  On the reduced residue classes modulo $ms$ lying over
$a\pmod m$, the function $(b/s)$ is a nontrivial character of the
kernel of reduction to $m$.  Its sum is zero.  If an exceptional
character is present, its contribution to this weighted sum also
vanishes unless its restriction to that kernel is exactly $(\cdot/s)$.
Call such a character aligned.  By QP2 at most one aligned index can
have $1-\beta<c_P/\log R$: two such characters must have the same
primitive inducer, and restricting to the kernel at any prime dividing
one of the two squarefree indices but not the other contradicts equality.
The Chinese remainder theorem provides the required test unit.
Take this index for $s_*$, or $s_*=0$ if none exists.

For every other index an aligned exceptional term is at most
$r\operatorname{Li}(U)$, with
$r\le4\log U\,e^{-c_P\log U/\log R}$.
The relative errors from (QP.9), summed over the fiber, are at most
$k$ times the sum of its positive main terms, where
$k=C(e^{-c\log U/\log(ms)}+e^{-c\sqrt{\log U}})$.
If alignment occurs the unweighted exceptional character sums to zero
on the fiber, so its total main term is the ordinary one; if it does
not occur the weighted main term is zero.  In either case division by
the total count costs at most $(1-k)^{-1}$, not $\varphi(s)$.
Removing primes below $\sqrt U$ changes the weighted sum by at most
$\sqrt U$; primes dividing $s$ have Jacobi value zero.

Uniformly for (QP.10),
$\log U/(\log m+A\log Y)\gg_{A,\epsilon}\log^2Y$.
Thus $r+k=o(Y^{-B})$, while (QP.11) gives
$\sqrt U/|P|\le U^{-2/5}=o(Y^{-B})$.
After increasing the single threshold these bounds prove (QP.12),
the literal \texttt{twistK} conclusion with arbitrary fixed $A,B$.

\paragraph{QP8: the coefficient estimates have absolute constants.}
For $y\ge3$ put $c_y(n)=\Lambda(n)\min(1,y/n)$ and, for $N\ge1$,
\[
 b_p=\sum_{p^j\le N}c_y(p^j),\quad
 R_m=\sum_{\substack{p^j\le N\cr j\ {\rm even}\ {\rm or}\ p\mid m}}c_y(p^j).
\]
The elementary binomial-coefficient proof of Chebyshev's bound gives
$\sum_{p\le t}\log p\le t\log4$ and $\psi(t)\le Ct$.
For $p\le y$, splitting the geometric series at
$\lfloor\log y/\log p\rfloor$ gives $b_p\le3\log y$.
For $p>y$, it gives $b_p\le2y\log p/p\le2\log y$.
The same Chebyshev bound implies $\pi(y)\le Cy/\log y$ by separating
the primes below $\sqrt y$.  Partial summation of
$\sum_{p>y}\log^2p/p^2$ then proves
\[
 \sup_p b_p\le C\log y,\qquad \sum_p b_p^2\le Cy\log y.\tag{QP.13}
\]
The even powers in $R_m$ are bounded by
$\sum_{n\le N}\Lambda(n)\min(1,y/n^2)\le C\sqrt y$,
using $\psi(t)\le Ct$ and the decreasing kernel.
If every prime divisor of $m$ that is at most $N$ is at most $w$, the
other powers contribute at most $Cw\log y$.  Hence
\[
 R_m\le C(\sqrt y+w\log y),\qquad
 \sum_{n\le N}c_y(n)\le Cy\log(2N).                 \tag{QP.14}
\]
Finally the identity
$\log(\lfloor z\rfloor!)=\sum_{n\le z}\Lambda(n)\lfloor\lfloor z\rfloor/n\rfloor$,
Stirling's elementary integral bounds, and $\psi(z)\le Cz$ show
\[
                    \sum_{n\le z}\Lambda(n)/n\le\log z+C.\tag{QP.15}
\]
Only its leading coefficient $1$, not the numerical value of $C$, is
needed below.  These estimates are uniform in the cutoff $N$.

\paragraph{QP9: expand the mean square without assuming multiplicativity.}
Let $P$ consist of primes $q>N^2$ in one reduced class modulo $m$,
where $8\mid m$.  Let $|\kappa(n)|\le1$; no multiplicativity is imposed.
Put
\[
 S_N(y,q,\kappa)=\sum_{n\le N}c_y(n)\kappa(n)(n/q).
\]
Quadratic reciprocity on this fixed class groups its terms as
$S_N=b_0+\sum_{p\nmid m}b'_p(q/p)$, where the latter indices are odd
primes, $|b_0|\le R_m$, and $|b'_p|\le b_p$.
Even powers and powers of primes dividing $m$ enter $b_0$; the signs
from reciprocity and the prime $2$ are fixed by the residue class.
Suppose all the squarefree correlations in (QP.12), up to $N^2$, are
bounded by $\eta|P|$ except possibly $s_*$.
The diagonal of the square expansion is bounded by
$(R_m^2+\sum_pb_p^2)|P|$.
An exceptional correlation can occur only at $s_*=p$, contributing
at most $2R_m\sup b_p$, or at $s_*=pp'$ for distinct primes,
contributing at most $2(\sup b_p)^2$.
All remaining cross terms total at most
$\eta(\sum_{n\le N}c_y(n))^2|P|$, by the triangle inequality.
Therefore
\[
 \sum_{q\in P}|S_N(y,q,\kappa)|^2\le |P|\left\{
 \sum_pb_p^2+R_m^2+2B^2+2BR_m+
                        \eta\left(\sum_{n\le N}c_y(n)\right)^2\right\},
 \quad B=\sup_pb_p.                                    \tag{QP.16}
\]
This is the source's \texttt{meanSq\_abs} inequality, obtained by an
explicit prime-power partition.  It explains why one exceptional
squarefree index does not destroy the desired mean square.

\paragraph{QP10: the actual polynomial-window mean square.}
Fix $0<\epsilon\le1$, $K\ge1$, set $Z=(Y^K)^{400}$ and
$N=\lfloor Z\rfloor$.  In addition to (QP.10), suppose every prime
dividing $m$ is at most $Y^{\epsilon/5}$ or exceeds $Z$.
QP7 with $A=800K$, $B=K+1$ gives the correlations for (QP.16), with
$\eta=Y^{-(K+1)}$.  For sufficiently large $Y$,
$\sqrt U>N^2$, so its prime-size condition also holds.
For $Y\le y\le Y^K$, take $w=Y^{\epsilon/5}$ in QP8.
Uniformly on this interval, $w\log y\le\sqrt y$ eventually.
Thus $R_m^2\le Cy$, and (QP.13) bounds the diagonal by $Cy\log y$.
The two exceptional terms are at most $C(\log^2y+\sqrt y\log y)$.
The final term divided by $y\log y$ is at most
$CY^{-1}(400K\log Y)^2/\log Y$, which tends to zero.
We have proved that one absolute $C>0$ satisfies
\[
 \forall y\in[Y,Y^K]\ \forall\kappa,quad |\kappa(n)|\le1
 \Longrightarrow\quad
 \sum_{q\in P}|S_N(y,q,\kappa)|^2\le Cy\log y\,|P|.   \tag{QP.17}
\]
The threshold depends on $\epsilon,K$, but $C$ does not.
This is the literal stronger \texttt{meanSqK} endpoint, including its
quantification over all bounded coefficient sequences.

\paragraph{QP11: a maximal estimate for a continuum of arguments.}
For fixed $q$, the function $S_N(y,q,\kappa)$ is continuous and piecewise
linear, with right derivative
\[
 D_q(v)=\sum_{v<n\le N}\Lambda(n)\kappa(n)(n/q)/n.
\]
The exact identity
$vD_q(v)=S_N(v,q,\kappa\,1_{n>v})$
allows (QP.17) to be applied to this derivative, because the truncated
sequence is still bounded by one.  Hence
$\sum_q|D_q(v)|^2\le C(\log v)|P|/v$.
On a block $[a,b]\subseteq[Y,Y^K]$ with $b\le2a$, the fundamental
theorem of calculus and Cauchy--Schwarz give, for every $y_q\in[a,b]$,
\[
 |S_N(y_q,q,\kappa)|^2
 \le2|S_N(a,q,\kappa)|^2+2a\int_a^b|D_q(v)|^2dv.
\]
Summing and integrating the derivative bound proves
\[
 \sum_{q\in P}|S_N(y_q,q,\kappa)|^2
                                  \le4Ca\log(2a)|P|.\tag{QP.18}
\]
All sums are finite and $D_q$ is a step function; thus no measurability
or exchange-of-limit issue occurs.  In particular each prime may choose
its own witness $y_q$.  This is the step that upgrades pointwise control
to control throughout a real interval.

\paragraph{QP12: discard only a small fraction of primes.}
Let $\tau>0$.  If a prime is bad on $[a,b]\subseteq[a,2a]$, choose a
witness with $|S_N(y_q,q,\kappa)|>y_q^{1/2+\tau}/2$.
By (QP.18), the proportion of such primes is at most
$16C\log(2a)a^{-2\tau}\le(32C/\tau)a^{-\tau}$ for $a\ge2$.
Use blocks $[2^iY,\min(2^{i+1}Y,Y^K)]$.
Every bad prime belongs to at least one block, so the geometric series
gives
\[
 \frac{\#\{q\in P:\ \exists y\in[Y,Y^K],\
                    |S_N(y,q,\kappa)|>y^{1/2+\tau}/2\}}{|P|}
 \le \frac{32C}{\tau(1-2^{-\tau})}Y^{-\tau}.          \tag{QP.19}
\]
This tends to zero uniformly in $P,N,\kappa$ whenever (QP.17) holds.
For a fixed finite family of at most $R$ coefficient sequences, choose
the threshold so that the right side is at most $\eta/(2(R+1))$ and
take their union.  This is the precise \texttt{bad\_share} supplier.

\paragraph{QP13: extend the explicit formula to the needed cutoff heights.}
Suppose $\chi$ is primitive nonprincipal, its conductor is at most $Q$,
and it has no zero with $\Re s>399/400$, $|\Im s|\le Q$.
Reflection by CC1 gives $1/400\le\Re\rho\le399/400$ for its nontrivial
zeros below height $Q$ (apply conjugation as well as reflection).
Take $T=Q/2$, with $Q$ sufficiently large.
The contour proof in CC8 also works when $T>u$, provided one retains
the left-edge term $Cu^{-1/2}\log^2(N(T+2))$ instead of absorbing it.
To check this extension, the Perron kernel bounds the term nearest $u$
by $C\log(2u)$; the remaining terms sum to
$Cu\log^2(2u)/T$ irrespective of the order of $u$ and $T$.
Choose the good horizontal edges in $[T,T+1]\subseteq[0,Q]$ as in CC8.
Their cost is $Cu\log^2(Qu)/Q$, and the left edge is the term just stated.
The pole-at-zero constant is $O(\log^2Q)$ by the reflected lower bound
and CC3.  The reciprocal zero weights sum to $C\log^2Q$ by CC7,
now with the fixed lower bound $1/400$ at small height.
It follows that for every $u\ge3$,
\[
 |\psi_\chi(u)|\le C\{u^{399/400}\log^2Q
       +(u/Q)\log^2(Qu)+\log(Qu)+\log^2Q
                              +u^{-1/2}\log^2Q\}.     \tag{QP.20}
\]
This inequality concerns this zero-free box only; it asserts no
zero-free strip beyond height $Q$.

\paragraph{QP14: a two-range tail estimate supplies polynomial truncation.}
Put $\lambda=\log Q$, and let $V=\exp(B\lambda^2)$, where $B$ is a
sufficiently large fixed absolute constant.  The real-zero exclusion
required in CC10 holds with $c_0=1/400$: real zeros have imaginary part
zero and are therefore in the box at every height.
CC11--CC12 justify convergence and partial summation of
\[
 T_\chi(Z):=-L'/L(1,\chi)-\sum_{n\le Z}\Lambda(n)\chi(n)/n
       =-\psi_\chi(Z)/Z+\int_Z^\infty\psi_\chi(u)/u^2\,du.\tag{QP.21}
\]
For $3\le Z\le V$, integrate (QP.20) up to $V$.
The zero term contributes at most $C Z^{-1/400}\lambda^2$;
the $1/Q$ term contributes at most $C\lambda^6/Q$;
the remaining terms contribute at most
$CZ^{-1}\log^2(QZ)$, including the endpoint in (QP.21).
Above $V$, CC10 gives an integral bounded by
\[
 C\int_{B\lambda^2}^{\infty}(v+\lambda)^2
             e^{-cv/(\lambda+\sqrt v)}dv.
\]
For $v\ge B\lambda^2$, $v/(\lambda+\sqrt v)\ge\sqrt v/2$.
Substitute $r=\sqrt v$ and integrate a fixed polynomial times
$e^{-cr/2}$.  Choose $B$ so large that this is at most $C\lambda^6/Q^2$.
We obtain the useful uniform estimate
\[
 |T_\chi(Z)|\le C\{Z^{-1/400}\lambda^2+
                  \lambda^6/Q+Z^{-1}\log^2(QZ)\}.     \tag{QP.22}
\]
This estimate extends to a nonprincipal character $\chi$ at any level
$N\le Q$.  Its primitive inducer $\chi^*$ is nonprincipal, and the deleted
Euler factors are nonzero on $\Re s>0$, so $\chi^*$ has the same required
zero-free box.  The difference of the two tails is a sum over the finitely
many deleted primes.  At each such prime,
\[
 \sum_{p^j>Z}\frac{\log p}{p^j}\le\frac{2\log p}{Z}.
\]
Summing gives $|T_\chi(Z)-T_{\chi^*}(Z)|\le2\log N/Z$, already
absorbed in (QP.22).  Thus (QP.22) holds under the literal
\texttt{GoodC} hypothesis, which does not require primitivity.

Let $x=Y^K$, $Z=x^{400}$ and
$Y^{\epsilon/6}\le\lambda\le Y^{\epsilon/4}$, where
$0<\epsilon\le1$, $K\ge1$.  Then $3\le Z\le V$ eventually and
\[
 xZ^{-1/400}\lambda^2=\lambda^2\le Y^{\epsilon/2},
 \qquad x\lambda^6/Q\longrightarrow0,
 \qquad xZ^{-1}\log^2(QZ)\longrightarrow0.
\]
The last two limits are uniform in the stated range of $\lambda$.
Thus $x|T_\chi(Z)|\le Y^\epsilon$ eventually, exactly \texttt{tailK}.
For $y\le x$, the exact centered identity is
\[
 G_\chi(y)=-\sum_{n\le Z}\Lambda(n)\chi(n)\min(1,y/n)
                                                 -yT_\chi(Z).\tag{QP.23}
\]
It follows by expanding the finite sum and (QP.21), with the same
$+yL'/L(1,\chi)$ centering as CC12.  Taking $K=1$ and using the
coarser form of (QP.15) also supplies \texttt{kappaC}, with an absolute
coefficient of $\log Y$.

\paragraph{QP15: the exact coefficient four follows from a log-log bound.}
In (QP.22) instead choose $Z=\lambda^{1000}$.
The first term is $\lambda^{-1/2}$, and the other two are bounded
absolutely for sufficiently large $\lambda$.
By (QP.15),
\[
 |L'/L(1,\chi)|\le1000\log\log Q+C.                  \tag{QP.24}
\]
This proves the source's \texttt{kappa4\_loglog} conclusion by the
two-range explicit formula, without the longer differenced-contour route.
If $0<\epsilon<1/52$, $A>0$, and
$Y^{\epsilon/6}\le\log Q\le A Y^{\epsilon/5}\log^2Y$, then
\[
 1000\log\log Q+C
 \le200\epsilon\log Y+2000\log\log Y+1000\log A+C
                                                    \le4\log Y
                                                               \tag{QP.25}
\]
for all sufficiently large $Y$.  The strict gap is
$4-200/52=2/13>0$.  Thresholds may depend on $\epsilon,A$, but not on
$Q,N,\chi$.  This proves the literal \texttt{kappa4} bound needed in
\texttt{PairGood}; a bound with an unspecified larger coefficient would
not have sufficed.

\paragraph{QP16: a hybrid density theorem counts bad primitive boxes.}
Apply (HD) at $\sigma=399/400$, $T=Q$, $h=1/100$.
The exponent is
$3(2/400+1/100)=9/200<1/4$.  Each primitive character of conductor
at most $Q$ whose box fails has at least one zero in this rectangle,
so the number of such nonprincipal characters is at most
\[
                             C Q^{9/200}\le C Q^{1/4}.\tag{QP.26}
\]
The one primitive principal character may be added separately.  In
particular a totalized value assigned to its pole at one is not counted
as an analytic zero in (HD); it costs at most one character and is
absorbed by changing $C$.  This is the existential
\texttt{ZeroDensityBox} endpoint actually used
by \texttt{count\_bad\_core}.  No explicit enormous numerical constant
from another density interface is claimed.
Moreover, if each bad prime $q>F$ is assigned a primitive character of
conductor $fq\le Q$ with $f\le F$, this assignment is injective.
Indeed $fq=f'q'$ implies $q\mid f'$ or $q=q'$; the first case is
impossible since $0<f'\le F<q$.  Thus the same $CQ^{1/4}$ bound applies
to the number of bad primes.  The argument permits the character to
depend on the prime in any manner.

\paragraph{QP17: primitive twists and induced-character quantifiers.}
Let $\xi$ be primitive of conductor $f$, and let $q$ be an odd prime
with $q>f$.  The quadratic character $(\cdot/q)$ is primitive modulo
$q$.  The Chinese remainder decomposition of unit groups shows that
$\xi(\cdot)(\cdot/q)$ is primitive of conductor $fq$: factoring through
a proper divisor would factor one of the two primitive components.
It is nonprincipal because its $q$ component is nontrivial.

There is also a useful uniqueness assertion in the source's
\texttt{Ind} convention.  If two primitive characters agree on all
integers coprime to a positive ambient modulus $M$, then they have the same
conductor and are the same primitive character.  To prove this, compare
them at a unit modulo the product of their conductors and use the
Chinese remainder theorem to choose a congruent representative avoiding
the primes of $M$ not in that product.  They induce the same character
at the product level; uniqueness of the conductor proves the assertion.
In the uses here the prescribed twist is a character at a known common
level, so this argument identifies every primitive character quantified
by \texttt{Ind}, not just a preferred choice of representative.
Removing an unused extra prime from the ambient level is justified by
the same argument.  This supplies the \texttt{ind\_drop} step when the
second quadratic exponent is zero.

\paragraph{QP18: the generic good-prime stage.}
Fix $0<\epsilon\le1/4$, $\tau>0$, $K\ge1$, $\eta>0$, and an integer
$R\ge0$.  In admitted data satisfying the prime-divisor condition of
QP10, let $1\le F\le\sqrt U$ and let $\Xi$ be any family of at most
$R$ primitive characters, each of conductor at most $F$.
For $q\in P$, QP17 gives nonprincipal primitive twists of conductor at
most $FU$; QP16 shows that at most $C(FU)^{1/4}\le CU^{3/8}$ primes
can have a bad box for one of these twists.
One can assign a single witness to each such prime, so this count does
not acquire a factor $R$.
By (QP.11), this is less than $\eta|P|/2$ eventually.

For each $\xi\in\Xi$, use QP12 with $\kappa(n)=\xi(n)$.
Outside another set of at most $\eta|P|/2$ primes, all its truncated
sums satisfy $|S_N(y,q,\xi)|\le y^{1/2+\tau}/2$ throughout
$[Y,Y^K]$.  Here $N=\lfloor(Y^K)^{400}\rfloor$.
For $Q=FU$, (QP.10) and $F\le\sqrt U$ give
$Y^{\epsilon/6}\le\log Q\le Y^{\epsilon/4}$ eventually.
QP14 and $Y^\epsilon\le y^{1/2+\tau}/2$ for $\epsilon\le1/4$
then give $|G_\chi(y)|\le y^{1/2+\tau}$.
QP14's \texttt{kappaC} supplies one absolute logarithmic-derivative
coefficient, and QP17 transfers these conclusions to every primitive
inducer at every permitted ambient level.
Thus all but $\eta|P|$ primes satisfy the full \texttt{BoxAll} conditions
for every base character.  For $\eta<1$ at least one such prime exists.
The threshold depends only on the fixed displayed parameters, including
$R$, and is uniform in the moving family and in $F$.

\paragraph{QP19: the first prime can be chosen uniformly over all base characters.}
Fix $N_0>0$, $0<\epsilon<1/52$, $\tau>0$, $K\ge1$.
Set $z=Y^{\epsilon/5}$ and
$P_1=\{q:\sqrt{e^z}<q\le e^z,\ q\equiv3\pmod8,\ q\text{ prime}\}$.
Use QP18 with $m=8$, $a=3$, $U=e^z$, $F=N_0$, and the primitive
inducers of all characters modulo $N_0$.
Their number is at most $\varphi(N_0)$; $F\le\sqrt U$ eventually.
The admitted inequalities and the prime-divisor condition hold for
large $Y$.  Thus $P_1$ is nonempty and all but any prescribed small
fraction satisfy \texttt{GoodQ} for every character modulo $N_0$.
Also $q>N_0$ and $q>(Y^K)^{400}$ uniformly for these primes eventually,
because $\log q>z/2$ dominates every fixed multiple of $\log Y$.
The constant in the \texttt{Box} records is independent of $N_0$.

\paragraph{QP20: a second progression with one large modulus prime.}
For any $q\in P_1$ put
$M_o=q\prod_{3\le p\le z, p\ {\rm prime}}p$.
The conditions $d\equiv7\pmod{16}$ and $d\equiv-1\pmod{M_o}$
give a single reduced class $a_0$ modulo $m_0=16M_o$ by CRT.
Let $U=e^{z'}$, where
\[
 z\log^2Y\le z'\le K^2 z\log^2Y.                    \tag{QP.27}
\]
The elementary primorial bound gives
$\log m_0\le\log16+z+z\log4\le4z$ eventually.
Apply QP18 with the enlarged exponent $\epsilon'=6\epsilon/5$.
Its lower admitted bound is $Y^{\epsilon'/6}=z$;
its upper one follows from $K^2\le Y^{\epsilon/25}$ eventually.
The remaining admitted inequality follows from $\log m_0\le4z$.
Every prime divisor of $m_0$ is at most $Y^{\epsilon'/5}$ except $q$,
which is larger than $(Y^K)^{400}$ by QP19.  This is precisely why the
large-prime alternative in QP10 was retained.

Take $F=N_0e^z$ and the primitive inducers of
$\xi(\cdot)(\cdot/q)^i$, $\xi\bmod N_0$, $i=0,1$.
There are at most $2\varphi(N_0)$ of them, their conductors are at most
$N_0q\le F$, and $F\le\sqrt U$ by (QP.27) for large $Y$.
Thus the second prime set is nonempty and almost all its elements
satisfy \texttt{GoodD} simultaneously for these base characters.
All thresholds are uniform over $q\in P_1$ and $z'$ in (QP.27).
The congruences also imply $q<d$ and that every odd prime $p\le z$,
and $p=q$, divides $d+1$.

\paragraph{QP21: positive proportions give the actual simultaneous pair.}
Take the good proportion at both stages to be at least $1/2$.
QP19 supplies a nonempty good subset of $P_1$; for every prime in that
subset, QP20 supplies a nonempty good subset of its second progression.
Choose $q$ and then $d$ from these subsets.  Use the maximum of the two
absolute \texttt{Box} constants, and the maximum of their finitely many
thresholds.  The \texttt{Box} conditions are monotone in that constant
and in a decrease of the upper endpoint $X$.
We obtain one absolute $C>0$ such that, for every fixed
$\epsilon,\tau,K,N_0$ as above and all sufficiently large $Y$, every
$X\le Y^K$ and every $z'$ in (QP.27) have primes $q,d$ satisfying
\texttt{Rows}$(Y^{\epsilon/5},z',q,d)$, $N_0<q$,
\texttt{GoodQ} and \texttt{GoodD}.  In particular both windows,
both prime congruences, their ordering, and every divisibility
condition hold for the same pair.
This proves the literal \texttt{pair\_share} and \texttt{pair\_exists}
conclusions under the named classical inputs above; positivity of the
proportions is used together with positivity of the actual prime counts.

\paragraph{QP22: install the complete PairGood record.}
Given $M>0$, $0<\epsilon<1/52$, $A\ge1$, apply QP21 with
$N_0=8M$, $\tau=\epsilon/2$, $K=A$, $X=Y^A$, and
$z'=Y^{\epsilon/5}\log^2Y$.
The resulting primes are coprime to $M$ and satisfy the literal
\texttt{WinPair} congruences and logarithmic windows.
For a character $\xi\bmod M$ and $(i,j)\in\{0,1\}^2\setminus\{(0,0)\}$,
inflate $\xi$ to modulus $8M$ and consider
$\xi(n)(n/q)^i(n/d)^j$ on integers coprime to $2Mqd$.
If $j=0$ then $i=1$; QP17 removes the unused $d$ and \texttt{GoodQ}
applies.  If $j=1$, \texttt{GoodD} applies for either value of $i$.
Every primitive inducer is nonprincipal and is covered by this argument.

Its height parameter is exactly
\[
 Q_j=8M\exp(Y^{\epsilon/5})
            \exp(Y^{\epsilon/5}\log^2Y)^j,\qquad
 Y^{\epsilon/6}\le\log Q_j\le3Y^{\epsilon/5}\log^2Y.  \tag{QP.28}
\]
Thus (QP.25) gives $|L'/L(1,\chi)|\le4\log Y\le4\log y$.
The centered bound is $|G_\chi(y)|\le y^{1/2+\epsilon/2}$
for every real $Y\le y\le Y^A$; on the finite sum $n\le y$,
$\max(y/n-1,0)=y/n-1$, so this is exactly \texttt{An.gChi}.
Applying the same argument to every real primitive base character with
conductor dividing $8M$ gives the required \texttt{ZeroBox} at height
$Q_j$.  Therefore the same $q,d$ satisfy all four clauses of
\texttt{PairGood}, with its universal induced-character quantifiers.
This proves \texttt{UpperTheory.Top.PairEx}.

\paragraph{QP23: real-character factorization gives the actual host exclusion.}
Let an already specified host $H$ have distinguished primes $q,d$ from
QP22, with $(H.QS,qd)=1$, every prime of $H.QS$ dividing $2M$, and
$\operatorname{ordCompl}_2(H.QS)\mid M$.
For a real primitive character belonging to $H$, its conductor divides
$H.QS\,qd$.  CRT decomposes it into its prime-power components.
At an odd prime the subgroup $1+p\mathbb Z/p^r\mathbb Z$ has odd
order, so a quadratic character is trivial there; its conductor
exponent is at most one.  At $2$, every unit congruent to $1\pmod8$
is a square (lift its odd square root one binary digit at a time), so
the conductor exponent is at most three.
The $q$ and $d$ components are therefore either trivial or the unique
quadratic characters modulo those primes.  Removing them leaves a
real primitive character $\rho$ with
$\rho.N\mid8\operatorname{ordCompl}_2(H.QS)$ and
\[
 \chi.N=\rho.N\,q^i d^j,\qquad
 \chi(n)=\rho(n)(n/q)^i(n/d)^j\quad((n,\chi.N)=1).    \tag{QP.29}
\]
If $q\mid\chi.N$ or $d\mid\chi.N$, then $(i,j)\ne(0,0)$.
QP22 gives the box for $\chi$; its height is positive and hence includes
every real zero.  This proves \texttt{RealBox H}:
$L(\beta,\chi)\ne0$ for every $\beta>399/400$ of the specified real
characters.  Moreover (QP.29) and the upper prime windows give the
literal conductor bound
\[
 \chi.N\le8\operatorname{ordCompl}_2(H.QS)
          e^{Y^{\epsilon/5}}
          e^{Y^{\epsilon/5}\log^2Y\,1_{d\mid\chi.N}}. \tag{QP.30}
\]
This is the same-host \texttt{realBox} endpoint; no replacement host
or enlargement of its prescribed fields has occurred.

\paragraph{QP24: close the exceptional-zero premise in the thin host estimate.}
Apply CC19--CC20 to the very host in QP23.
Whenever its other displayed hypotheses hold,
$\operatorname{ordCompl}_2(H.QS)\le B$ and
$\log H.\mathrm{condN}\le2x^{\epsilon/5}\log^2x$,
QP23 supplies the \texttt{RealBox H} premise previously retained there.
The conclusion is the actual \texttt{H.ChebyshevCharacterBound}
with the same centering, the same lower range
$y\ge\frac13\exp(x^{\epsilon/3})$, and one uniform positive decay
constant.  Thus the CC thin estimate now has a concrete prime-choice
supplier for its exceptional-zero exclusion whenever the chosen host
uses the pair of QP22 and satisfies QP23's level conditions.
This assertion does not construct an arbitrary host, prove a missing
geometric window, or supply the global automorphic input.  Those
requirements are independent of the real-character factorization.

\paragraph{QP25: the remaining window premise is not lost in a composition.}
The source's \texttt{winPairs} conclusion follows from QP22 and an
independently supplied \texttt{VWinAt}$(\psi,401)$, with the exact order
of quantifiers preserved.  To check the rescaling, choose the integer
page function with
$u^{2+8\epsilon}\le2^{e(u)}\le2u^{2+8\epsilon}$.
Let $J\ge1$ and its initial threshold come from \texttt{VWinAt}.
For a desired factor $c\ge1$, put $k=c^{1/(2+8\epsilon)}$.
Choose $M$ above the PairEx threshold and above $k^3,k^{400}$, and take
$X\ge M^J$.  QP22 gives a good pair at $Y=X^{1/J}$.
The window premise supplies $W\ge Y$, $W^2\le X$ and its geometric
window for that same pair.  Put $Z=kW$.
For $Z\le x\le Z^{3/2}$, one has
$W\le x/k\le x\le W^2$, because $W\ge k^3$.
Use the page $e(x/k)$ and multiply the preceding page inequalities by
$c$.  The geometric lower cutoff transfers because
$k^{400}\le x/k$, which gives
$x^{400}\le(x/k)^{401}$; its upper cutoff is monotone.
PairGood transfers from $X$ to $x$ by monotonicity.
This proves the literal \texttt{WinDatum} conclusion on
$[Z,Z^{3/2}]$ for arbitrarily large $Z$, conditional on the original
\texttt{VWinAt} input.  It neither proves that input nor advances any
global RH/GRH branch solely by reusing this implication.
